% Options for packages loaded elsewhere
\PassOptionsToPackage{unicode}{hyperref}
\PassOptionsToPackage{hyphens}{url}
\PassOptionsToPackage{dvipsnames,svgnames,x11names}{xcolor}
\documentclass[
  11pt,
]{article}
\usepackage{xcolor}
\usepackage[margin=25mm]{geometry}
\usepackage{amsmath,amssymb}
\setcounter{secnumdepth}{-\maxdimen} % remove section numbering
\usepackage{iftex}
\ifPDFTeX
  \usepackage[T1]{fontenc}
  \usepackage[utf8]{inputenc}
  \usepackage{textcomp} % provide euro and other symbols
\else % if luatex or xetex
  \usepackage{unicode-math} % this also loads fontspec
  \defaultfontfeatures{Scale=MatchLowercase}
  \defaultfontfeatures[\rmfamily]{Ligatures=TeX,Scale=1}
\fi
\usepackage{lmodern}
\ifPDFTeX\else
  % xetex/luatex font selection
\fi
% Use upquote if available, for straight quotes in verbatim environments
\IfFileExists{upquote.sty}{\usepackage{upquote}}{}
\IfFileExists{microtype.sty}{% use microtype if available
  \usepackage[]{microtype}
  \UseMicrotypeSet[protrusion]{basicmath} % disable protrusion for tt fonts
}{}
\makeatletter
\@ifundefined{KOMAClassName}{% if non-KOMA class
  \IfFileExists{parskip.sty}{%
    \usepackage{parskip}
  }{% else
    \setlength{\parindent}{0pt}
    \setlength{\parskip}{6pt plus 2pt minus 1pt}}
}{% if KOMA class
  \KOMAoptions{parskip=half}}
\makeatother
\usepackage{listings}
\newcommand{\passthrough}[1]{#1}
\lstset{defaultdialect=[5.3]Lua}
\lstset{defaultdialect=[x86masm]Assembler}
\usepackage{longtable,booktabs,array}
\usepackage{calc} % for calculating minipage widths
% Correct order of tables after \paragraph or \subparagraph
\usepackage{etoolbox}
\makeatletter
\patchcmd\longtable{\par}{\if@noskipsec\mbox{}\fi\par}{}{}
\makeatother
% Allow footnotes in longtable head/foot
\IfFileExists{footnotehyper.sty}{\usepackage{footnotehyper}}{\usepackage{footnote}}
\makesavenoteenv{longtable}
\setlength{\emergencystretch}{3em} % prevent overfull lines
\providecommand{\tightlist}{%
  \setlength{\itemsep}{0pt}\setlength{\parskip}{0pt}}
\emergencystretch=3em
\AtBeginDocument{\Urlmuskip=0mu plus 2mu}
\lstset{breaklines=true,basicstyle=\small\ttfamily,columns=fullflexible,keepspaces=true}
\setlength{\tabcolsep}{4pt}
\usepackage{iftex}
\ifXeTeX
  \usepackage{newunicodechar}
  \makeatletter
  \newcommand{\DeclareUnicodeCharacter}[2]{%
    \begingroup
    \uccode`\~="#1\relax
    \uppercase{\def\temp{~}}%
    \expandafter\endgroup
    \expandafter\newunicodechar\expandafter{\temp}{#2}%
  }
  \makeatother
\fi
\ifdefined\DeclareUnicodeCharacter
\DeclareUnicodeCharacter{2192}{\ensuremath{\to}}
\DeclareUnicodeCharacter{2190}{\ensuremath{\leftarrow}}
\DeclareUnicodeCharacter{2194}{\ensuremath{\leftrightarrow}}
\DeclareUnicodeCharacter{21D2}{\ensuremath{\Rightarrow}}
\DeclareUnicodeCharacter{21A6}{\ensuremath{\mapsto}}
\DeclareUnicodeCharacter{2212}{\ensuremath{-}}
\DeclareUnicodeCharacter{2264}{\ensuremath{\leq}}
\DeclareUnicodeCharacter{2265}{\ensuremath{\geq}}
\DeclareUnicodeCharacter{2260}{\ensuremath{\neq}}
\DeclareUnicodeCharacter{2113}{\ensuremath{\ell}}
\DeclareUnicodeCharacter{03A3}{\ensuremath{\Sigma}}
\DeclareUnicodeCharacter{03A6}{\ensuremath{\Phi}}
\DeclareUnicodeCharacter{0393}{\ensuremath{\Gamma}}
\DeclareUnicodeCharacter{221E}{\ensuremath{\infty}}
\DeclareUnicodeCharacter{2208}{\ensuremath{\in}}
\DeclareUnicodeCharacter{2209}{\ensuremath{\notin}}
\DeclareUnicodeCharacter{211A}{\ensuremath{\mathbb{Q}}}
\DeclareUnicodeCharacter{211D}{\ensuremath{\mathbb{R}}}
\DeclareUnicodeCharacter{2124}{\ensuremath{\mathbb{Z}}}
\DeclareUnicodeCharacter{2115}{\ensuremath{\mathbb{N}}}
\DeclareUnicodeCharacter{2200}{\ensuremath{\forall}}
\DeclareUnicodeCharacter{2203}{\ensuremath{\exists}}
\DeclareUnicodeCharacter{2286}{\ensuremath{\subseteq}}
\DeclareUnicodeCharacter{00D7}{\ensuremath{\times}}
\DeclareUnicodeCharacter{03B4}{\ensuremath{\delta}}
\DeclareUnicodeCharacter{2202}{\ensuremath{\partial}}
\DeclareUnicodeCharacter{03B2}{\ensuremath{\beta}}
\DeclareUnicodeCharacter{03B1}{\ensuremath{\alpha}}
\DeclareUnicodeCharacter{03C6}{\ensuremath{\phi}}
\DeclareUnicodeCharacter{03B3}{\ensuremath{\gamma}}
\DeclareUnicodeCharacter{00B9}{\ensuremath{^1}}
\DeclareUnicodeCharacter{00B2}{\ensuremath{^2}}
\DeclareUnicodeCharacter{2070}{\ensuremath{^0}}
\DeclareUnicodeCharacter{2080}{\ensuremath{_0}}
\DeclareUnicodeCharacter{2081}{\ensuremath{_1}}
\DeclareUnicodeCharacter{2082}{\ensuremath{_2}}
\fi
\usepackage{bookmark}
\IfFileExists{xurl.sty}{\usepackage{xurl}}{} % add URL line breaks if available
\urlstyle{same}
\hypersetup{
  pdftitle={PFP IV: Native Verification and the Custody of Finite Evidence},
  pdfauthor={Jeremy H. Carroll},
  colorlinks=true,
  linkcolor={Maroon},
  filecolor={Maroon},
  citecolor={Blue},
  urlcolor={Blue},
  pdfcreator={LaTeX via pandoc}}

\title{PFP IV: Native Verification and the Custody of Finite Evidence}
\author{Jeremy H. Carroll}
\date{2026-10-02 · Version 0.1.0 · Systems methods paper · Author-review
revision}

\begin{document}
\makeatletter
\begin{center}
{\LARGE \@title\par}
\vspace{0.8em}
{\normalsize \@author\par}
\vspace{0.5em}
{\small \@date\par}
\end{center}
\makeatother

\textbf{Code baseline:}
\nolinkurl{52cb5457f6c485ba548cfd292f6012a941ed22cf}

\subsection{Abstract}\label{abstract}

The Proposal Fidelity Protocol (PFP), formerly the Symbol Game,
separates untrusted proposal generation from checking a declared finite
property. This paper describes the concrete native verification
architecture at source commit
\nolinkurl{52cb5457f6c485ba548cfd292f6012a941ed22cf}. Its contribution
is an account of the interfaces and evidence boundaries required to
preserve that separation across Python objects, canonical encodings, a
Rust process, exact solver replay, and portable outcome records. A
worked four-vertex fixture illustrates the complete path: the host
reconstructs a carrier, recomputes packet and request identities,
replays an exact solve, compares the proposed normal form, and returns
exact laboratory evidence. The same result refuses admission because
three conformance laws remain unchecked and its attestation is unsigned.
We distinguish this restricted host route from the crate's broader
runtime dispatcher and additive foundation library. We also distinguish
the repository's tagged encoding and sorted JSON profiles from RFC 8785.
The system offers finite integrity and replay checks under explicit
assumptions; it does not establish producer authenticity, arbitrary
Python/Rust equivalence, a reproducible source-to-binary build, or
authority to deploy a result.

\subsection{1. The problem and the
contribution}\label{the-problem-and-the-contribution}

A generated answer may be syntactically acceptable, mathematically valid
on one finite carrier, and unsuitable for operational admission.
Collapsing those judgments into a single success flag loses essential
information. The risk grows when a Python caller serializes an object,
invokes a native solver, and later presents only an unqualified
``verified'' label.

PFP addresses this engineering problem by retaining distinct artifacts.
A proposal identifies what should be evaluated. The host reconstructs
what it can evaluate. An execution outcome records which checks
occurred, which exact values were obtained, and which obligations
remain. A policy decision determines whether the result may cross a
consequential boundary. The examined host's execution records confer no
operational authorization.

This paper contributes a source-pinned interface inventory, a complete
fixture trace, and a threat model for interpreting the emitted
commitments. It also identifies deliberately unavailable paths and
uneven canonicalization contracts. It is a systems and methods report.
Its bounded test cases are evidence about their exercised paths; a
passing test suite would not prove every reachable program behavior.

The scope is the
\href{https://github.com/dbccpage/the_symbol_game/tree/52cb5457f6c485ba548cfd292f6012a941ed22cf/omega_engine/crates_rs/omega_symbol_game}{Rust
\passthrough{\lstinline!omega\_symbol\_game!} crate}, its
\href{https://github.com/dbccpage/the_symbol_game/blob/52cb5457f6c485ba548cfd292f6012a941ed22cf/omega_engine/crates_rs/omega_symbol_game/src/protocol.rs}{protocol},
the Python
\href{https://github.com/dbccpage/the_symbol_game/blob/52cb5457f6c485ba548cfd292f6012a941ed22cf/omega_engine_symbol_game/ffi.py}{native
seam wrapper}, and the named transport and regression sources below. The
repository that holds those files is private. This deposit includes an
immutable copy of the audited source at the code baseline,
\nolinkurl{audited_source/pfp-iv-native-source-52cb5457.tar.gz}. The GitHub
links name that revision. The archive is the copy a reader can inspect.
The four-vertex calculation used here is restated in Section 6 and in
the companion study guide. PFP I is a separate preprint in this series
and is not required to check this fixture. The present paper explains
how an implementation records a finite calculation without enlarging its
meaning.

\subsection{2. Related work and design
assumptions}\label{related-work-and-design-assumptions}

Proof-carrying code studies a producer supplying evidence that a
consumer checks against a safety policy. The relevant architectural
precedent is explicit consumer-side checking. PFP's laboratory records
have a different scope and are not a proof that arbitrary generated code
satisfies a safety policy
\href{https://doi.org/10.1145/263699.263712}{(Necula, 1997)}.

Cryptographic comparison requires agreement about the bytes being
compared. RFC 8785 specifies JSON canonicalization rules covering
primitive serialization and deterministic property sorting. It also
requires preservation of string data rather than Unicode normalization.
The repository has multiple encoding profiles, so conformance cannot be
inferred from a function named \passthrough{\lstinline!canonical\_json!}
\href{https://doi.org/10.17487/RFC8785}{(Rundgren, Jordan, and Erdtman,
2020)}.

Our execution assumptions are explicit: the declared carrier is finite;
the arithmetic profile is known; the checker and native binary selected
by the launcher are the intended ones; the operating environment
executes those bytes; and any expected digest anchor is retained
independently of the artifact being checked. Invalidating one assumption
requires a narrower interpretation of the result. Exact arithmetic alone
does not discharge the other assumptions.

\subsection{3. Three interfaces, three different
capabilities}\label{three-interfaces-three-different-capabilities}

The crate exposes three distinguishable surfaces.

\begin{longtable}[]{@{}
  >{\raggedright\arraybackslash}p{(\linewidth - 4\tabcolsep) * \real{0.3333}}
  >{\raggedright\arraybackslash}p{(\linewidth - 4\tabcolsep) * \real{0.3333}}
  >{\raggedright\arraybackslash}p{(\linewidth - 4\tabcolsep) * \real{0.3333}}@{}}
\toprule\noalign{}
\begin{minipage}[b]{\linewidth}\raggedright
Surface
\end{minipage} & \begin{minipage}[b]{\linewidth}\raggedright
Entry
\end{minipage} & \begin{minipage}[b]{\linewidth}\raggedright
Evidence boundary
\end{minipage} \\
\midrule\noalign{}
\endhead
\bottomrule\noalign{}
\endlastfoot
Restricted prime host & Default executable operation,
\passthrough{\lstinline!evaluate\_request\_bytes!} & One same-carrier
candidate, exact replay, explicit admission result \\
Runtime dispatcher & \passthrough{\lstinline!invoke!},
\passthrough{\lstinline!--describe-runtime!},
\passthrough{\lstinline!dispatch\_runtime\_request!} & Per-action
payload and outcome contracts \\
Additive foundation & Rust library \passthrough{\lstinline!foundation!}
& Finite exhaustion, extension, packet, and composition witnesses \\
\end{longtable}

The
\href{https://github.com/dbccpage/the_symbol_game/blob/52cb5457f6c485ba548cfd292f6012a941ed22cf/omega_engine/crates_rs/omega_symbol_game/src/main.rs}{executable
source} also offers \passthrough{\lstinline!--describe!},
\passthrough{\lstinline!--derive-carrier-manifest!}, and
\passthrough{\lstinline!classify-board!}.
\passthrough{\lstinline!--describe!} reports the restricted host
profile; \passthrough{\lstinline!--describe-runtime!} reports adapter
names derived from the runtime action table. Descriptions are useful
discovery artifacts, but advertising an action does not establish that
every mathematical input or every external use is supported.

The
\href{https://github.com/dbccpage/the_symbol_game/blob/52cb5457f6c485ba548cfd292f6012a941ed22cf/omega_engine/crates_rs/omega_symbol_game/src/invoke.rs}{dispatcher}
includes functional evaluation and replay, exact matrix operations,
cochain checks, code-change actions, and other named adapters. A runtime
action called \passthrough{\lstinline!prime\_admit!} must be interpreted
under that action's own contract. Its existence does not remove the
restricted host's refused admission result. Likewise, a finite
conservative-extension witness in the foundation library does not wire a
carrier-changing lift into the default host protocol.

Python's \nolinkurl{omega_engine_symbol_game.ffi.RustSymbolGameHost}
delegates process location and invocation to the operating-system
adapter. The supported
\passthrough{\lstinline!OMEGA\_SYMBOL\_GAME\_HOST!} environment override
belongs to that seam. Convenience methods retain the full railway
outcome or typed refusal. Equivalent names in Python and Rust are
insufficient to establish functional parity: parity needs a declared
common input domain, output relation, and differential evidence.

A separate tutorial adapter,
\nolinkurl{omega_engine_providers.adapters.native_functional.compute_graph_phi1},
requires \nolinkurl{OMEGA_FUNCTIONAL_GRAPH_BIN} and
\nolinkurl{OMEGA_FUNCTIONAL_GRAPH_SHA256}. It validates the executable
digest and exact graph inputs, then parses a narrowly shaped native
result. This is a separate binary-selection contract. Configuring one
environment variable family does not configure the other.

\subsection{4. Packet identity and
canonicalization}\label{packet-identity-and-canonicalization}

The principal Python transport type is
\href{https://github.com/dbccpage/the_symbol_game/blob/52cb5457f6c485ba548cfd292f6012a941ed22cf/omega_engine_symbol_game/domain/entities/packet_envelope.py}{\passthrough{\lstinline!PacketEnvelopeV2!}}.
Its identity includes the task and packet family, carrier, initial
state, payload, constraints, protected boundaries, terminal target,
evidence ceiling, oracle policy, allowed transforms, budget, framing,
and source/parent bindings. The payload is frozen through the domain
canonical representation. Parsing the transport object does not create a
verified solve receipt or a continuation permit.

Three serialization distinctions must survive documentation:

\begin{longtable}[]{@{}
  >{\raggedright\arraybackslash}p{(\linewidth - 2\tabcolsep) * \real{0.5000}}
  >{\raggedright\arraybackslash}p{(\linewidth - 2\tabcolsep) * \real{0.5000}}@{}}
\toprule\noalign{}
\begin{minipage}[b]{\linewidth}\raggedright
Encoding surface
\end{minipage} & \begin{minipage}[b]{\linewidth}\raggedright
Actual contract at the audited revision
\end{minipage} \\
\midrule\noalign{}
\endhead
\bottomrule\noalign{}
\endlastfoot
Domain canonical algebra & Tagged values, including exact integers,
fractions, and supported compound types \\
\passthrough{\lstinline!PacketCodec!} helper & NFC normalization
followed by tagged domain encoding \\
Restricted host profile &
\passthrough{\lstinline!omega-sorted-json-no-float-v1!} for its outer
request structure \\
\end{longtable}

The
\href{https://github.com/dbccpage/the_symbol_game/blob/52cb5457f6c485ba548cfd292f6012a941ed22cf/omega_engine_symbol_game/domain/shared_kernel/canonical.py}{domain
canonical module} explicitly encodes type information. For example, an
integer and a string containing the same digits remain distinguishable.
It validates canonical compound encodings on reconstruction. The
\href{https://github.com/dbccpage/the_symbol_game/blob/52cb5457f6c485ba548cfd292f6012a941ed22cf/omega_engine_symbol_game/domain/packet_codec.py}{\passthrough{\lstinline!PacketCodec!}
helper} explicitly says its bytes are not RFC 8785. Its
\passthrough{\lstinline!verify!} method checks payload digest and
authority; it does not validate every field sealed into its
\passthrough{\lstinline!envelope\_sha256!}. It must not be substituted
for the richer protocol verifier when claiming complete envelope
custody.

The host's outer no-float profile and the Python tagged payload algebra
operate at different layers. The host tests include canonical tagged
float and fraction images, while theorem-bearing arithmetic uses exact
rational values. Accepting a typed representation in a payload does not
license floating-point approximation in an exact solver input.

Canonicalization establishes reproducible representation under a
profile. Hashing establishes an integrity comparison under an anchor.
Semantic equivalence requires an explicit projection or checker. None of
these steps, alone, authenticates who produced the record.

\subsection{5. Restricted host
execution}\label{restricted-host-execution}

The
\href{https://github.com/dbccpage/the_symbol_game/blob/52cb5457f6c485ba548cfd292f6012a941ed22cf/omega_engine/crates_rs/omega_symbol_game/src/host.rs}{host
implementation} admits the following narrow successful path. The request
uses \nolinkurl{omega.oesos.rust-prime-host-request.v1}, carries exactly
one candidate, and binds the host registry and build-stamped source
revision. The arithmetic domain is rational and the acceptance mode is
\nolinkurl{must_equal_replayed_normal_form-v1}. This restricted request
wire accepts weights and before/candidate coefficients only as minimal
signed 128-bit integer strings; it cannot express every rational input.
The solver computes with \passthrough{\lstinline!BigRational!}
internally. The broader \passthrough{\lstinline!invoke!} dispatcher has
its own input contract, which must be checked separately.

First, bounded intake rejects malformed, duplicate, unknown, or
oversized request data. The host recomputes request and packet
identities rather than trusting caller-supplied values. It validates the
carrier declaration, derives its manifest, and compares the version and
carrier hash to the supplied anchor. Before and candidate vectors must
match the edge dimension.

The host next performs the same-carrier identity checks and executes the
exact solve. The solve receipt remains inside the Rust process and is
replayed there. The candidate is accepted for this laboratory execution
only if it exactly equals the replayed rational normal form. A caller
cannot obtain that result by changing a candidate digest consistently.

Subsequent checks interpret exact final masses under an integral
\passthrough{\lstinline!u64!} policy; they do not round nonintegral
values to fit the policy. Conformance records preserve unchecked laws.
The supported route consumes one internal permit for an exact write and
another for terminalization. Those permit objects are not serialized;
the outcome carries only digest bindings for them.

Finally, an outcome carries rational summaries, typed refusal reasons,
commitments, and an outcome root. Even a fully executed route may return
refused admission. Exported bytes are a summary of the process evidence;
retaining the internal receipt across solve and replay prevents a Python
caller from manufacturing the particular Rust evidence object used in
that sequence.

\subsection{6. Worked example: problem change with refused
admission}\label{worked-example-problem-change-with-refused-admission}

The canonical
\href{https://github.com/dbccpage/the_symbol_game/blob/52cb5457f6c485ba548cfd292f6012a941ed22cf/omega_engine/crates_rs/omega_symbol_game/tests/host_contract.rs}{host
integration fixture} uses vertices 0, 1, 2, 3, unit-weight oriented
edges

\passthrough{\lstinline!(0,1), (1,2), (2,0), (1,3), (3,0)!}

and triangular faces \passthrough{\lstinline!(0,1,2)!},
\passthrough{\lstinline!(0,1,3)!}. Its before vector is
\passthrough{\lstinline!(3,2,2,1,1)!} and its candidate vector is
\passthrough{\lstinline!(0,0,0,0,0)!}. The fixture constructs the
manifest, computes the packet ID, obtains the host description, and
computes the request ID using exported helpers.

This carrier matters: deleting a face changes what closure means, even
if the edge list remains identical. Consequently the manifest commits to
faces as well as vertices and edges. A graph-only certificate cannot
silently become a certificate for this two-dimensional carrier.

The test specifies the following expected observations:

\begin{longtable}[]{@{}
  >{\raggedright\arraybackslash}p{(\linewidth - 2\tabcolsep) * \real{0.5000}}
  >{\raggedright\arraybackslash}p{(\linewidth - 2\tabcolsep) * \real{0.5000}}@{}}
\toprule\noalign{}
\begin{minipage}[b]{\linewidth}\raggedright
Observation
\end{minipage} & \begin{minipage}[b]{\linewidth}\raggedright
Expected value
\end{minipage} \\
\midrule\noalign{}
\endhead
\bottomrule\noalign{}
\endlastfoot
Execution status & \nolinkurl{EXECUTED_REPLAYED_EXACT_LAB_EVIDENCE} \\
Replayed candidate relation & Exact match \\
Intermediate gauge summaries & Phi1 = 7; Gamma2 = 12 \\
Final summaries & Phi1 = 0; Gamma2 = 0 \\
Conformance & \passthrough{\lstinline!Partial!}; three unchecked laws \\
Internal consumed letters & \passthrough{\lstinline!E,T!}; two budget
units \\
Attestation & \passthrough{\lstinline!unsigned\_lab!} \\
Admission & \passthrough{\lstinline!Refused!} \\
Admission reasons & \nolinkurl{PRIME000_CONFORMANCE_PARTIAL},
\passthrough{\lstinline!UNSIGNED\_LAB\_ATTESTATION!} \\
\end{longtable}

The intermediate values can be checked by hand: the before vector
circulates 3 + 2 + 2 = 7 around face \passthrough{\lstinline!(0,1,2)!}
and 3 + 1 + 1 = 5 around face \passthrough{\lstinline!(0,1,3)!}, so
Gamma2, the unweighted L1 norm of these face values, is 7 + 5 = 12,
while Phi1, the unit-weight L1 distance from the before vector to the
edge differences of vertex potentials, is 7, attained by the residual
\passthrough{\lstinline!(5,2,0,0,0)!} and matched from below by the unit
circulation around \passthrough{\lstinline!(0,1,2)!}. The two face sums are 7 and 5, so Gamma2 of the before vector is 12.
Face sums are unchanged by coboundaries. No potential takes
(3,2,2,1,1) to the zero vector. The jump is a problem change, not an
exact repair. The laboratory still reports final summaries 0 and 0 on
that replaced candidate, and admission stays refused. Gamma2 of the
before vector remains 12.

These are assertions of the test at the pinned source. Section 9
separates that source text from the retained historical native run. They
illustrate how numerical success and an unmet admission obligation
coexist. A consumer that displays only the final zeros would erase the
most important qualification.

The same suite supplies counterexamples. Changing a carrier while
retaining its old manifest anchor is refused, even when outer request
hashes are recomputed. Changing a candidate and recomputing its
surrounding identity still fails if it is not the replayed normal form.
Editing a field of a returned outcome breaks its root check; because the
root is unkeyed, recomputing it would restore only internal consistency,
not create a fresh native replay, and the outcome remains unsigned
laboratory evidence with refused admission. Each case targets a distinct
attempted shortcut through custody.

The novice interpretation is direct: the laboratory completed a
specified calculation, but the result lacks the additional evidence
needed for admission. ``Refused'' here is not a statement that
arithmetic failed.

\subsection{7. Resource boundaries and unavailable
workflows}\label{resource-boundaries-and-unavailable-workflows}

The restricted host accepts at most 1,048,576 request bytes. It bounds
declared vertices, edges, and faces, and the sum of vertices and faces,
at 4,096 each. Its preallocation fence also bounds four dense exact
matrices (d0, d1, and two augmented systems) at 1,048,576 cells each;
dimension products are checked before allocation. The supported
successful route needs two prime writes; the declared budget range is
2--64.

The timeout field is validated metadata (1 to 300,000 milliseconds), not
an in-process watchdog. A caller requiring a wall-clock kill boundary
must enforce it around the process and record that enforcement. Exact
rational operations can still be expensive within dimension limits; the
paper claims no general throughput or polynomial bound from these
fences.

Runtime \passthrough{\lstinline!invoke!} has its own 65,536-byte input
ceiling. The intake reader consumes at most the limit plus one witness
byte, which permits an oversized request to be rejected without reading
an unbounded stream. These limits belong to different interfaces and
should be reported separately.

The legacy Python \passthrough{\lstinline!run-turn!},
\passthrough{\lstinline!run-bout!}, and \passthrough{\lstinline!run!}
CLI routes are currently required to return
\passthrough{\lstinline!RUST\_FUNCTIONAL\_REQUIRED!} before creating
output.
\href{https://github.com/dbccpage/the_symbol_game/blob/52cb5457f6c485ba548cfd292f6012a941ed22cf/omega_engine_symbol_game/tests/test_functional_boundary.py}{Boundary
regressions} document that retirement. The supported offline catalog
commands and native \passthrough{\lstinline!gate11-self!} path provide
more precise entry points. In these engine paths, missing native
configuration is an unavailable calculation, not a zero defect or a
silently substituted Python estimate. The companion governor in this
package follows the same rule when a native executable is configured,
and without one it computes the minimum by exact enumeration (Section
10).

For the graph tutorial, the historical operator fixture supplies a small
hand check: claims (1,1,1) with penalties (1,2,2) have quotient minimum
3 and an attaining projected cochain (-2,1,1). The series' shared
foundation example uses penalties (1,2,3): its minimum is also 3, while
its zero-potential cost is 6 rather than this fixture's 5. A chosen
section's residual and the minimum over all sections are different
quantities. Comparing two Python wrappers that call the same function is
an API compatibility check, not independent solver agreement. The
weighted-cycle note, \nolinkurl{docs/completion_supplement.md}, gives
the argument and describes the companion code in this package.

\subsection{8. Threat model and evidence
claims}\label{threat-model-and-evidence-claims}

The adversary may control proposal text, serialized fields, stale
outcomes, and files offered as witnesses. The reviewed protections
target malformed intake, carrier drift, wrong candidate values, stale
commitments, and unsupported requests. The trusted computing base
includes the selected native executable, its dependencies, exact
arithmetic implementation, protocol code, and execution environment.

An unkeyed outcome root cannot authenticate the producer. An attacker
able to replace both a record and its expected anchor can recompute
both. The source revision is a build stamp; the host description does
not itself supply a trusted executable digest or signed
reproducible-build statement. Several source bindings are committed
names or digests whose underlying source bytes are not independently
reconstructed by the default host.

The interpretation must therefore remain conditional. An independently
pinned executable hash identifies executed bytes; it does not prove how
those bytes were built. A fresh replay supplies evidence about the
declared input and checker; it does not show that the carrier represents
an external physical system faithfully. Operational authorization
requires its own policy and trusted provenance.

This package also retains six small Lean modules from its former
manuscript in \passthrough{\lstinline!lean4/!}: incidence-sign and
signed column-sum lemmas, identity-path and quotient lemmas, rigidity of
one oriented two-point relation, the weighted null-space inequalities
and unique landing for one directed 3-cycle, a list-general
directional-derivative bound, and 60 elementary set and point-set
topology lemmas. They build with the pinned toolchain, and
\passthrough{\lstinline!lean4/axiom\_report.txt!} lists the axioms used
by each of the 86 theorems. None of them formalizes the native host, the
packet protocol, or a receipt field. No end-to-end refinement theorem
connecting every Python object, native action, receipt field, and Lean
statement is asserted here.

\subsection{9. Reproducibility and verification
plan}\label{reproducibility-and-verification-plan}

To repeat the native checks, extract
\nolinkurl{audited_source/pfp-iv-native-source-52cb5457.tar.gz}, record the
actual toolchain, and run only the relevant targets from the extracted
root. The principal native commands are:

\begin{lstlisting}
cargo test -j 1 --locked --manifest-path omega_engine/crates_rs/Cargo.toml -p omega_symbol_game --test host_contract
cargo test -j 1 --locked --manifest-path omega_engine/crates_rs/Cargo.toml -p omega_symbol_game --test foundation_contract
cargo test -j 1 --locked --manifest-path omega_engine/crates_rs/Cargo.toml -p omega_symbol_game --test invoke_contract
cargo test -j 1 --locked --manifest-path omega_engine/crates_rs/Cargo.toml -p omega_symbol_game --test native_cli_intake
\end{lstlisting}

\passthrough{\lstinline!host\_contract!} checks the worked path and
hostile variations. \passthrough{\lstinline!foundation\_contract!}
includes a cross-language golden packet vector, incomplete exhaustion
rejection, detector preservation, and composition/replay refusals. Those
shared vectors establish the declared examples, not parity on every
input. A historical author-run report records the four native targets on
2026-09-28 EDT (2026-09-29 UTC), using Rust 1.98.0 on Windows
(x86\_64-pc-windows-msvc). It reports
\passthrough{\lstinline!host\_contract!} 12 passed,
\passthrough{\lstinline!foundation\_contract!} 11,
\passthrough{\lstinline!invoke\_contract!} 32 and
\passthrough{\lstinline!native\_cli\_intake!} 9, with no failures. The
run commit is \nolinkurl{1c7092c3db2ab2fa65999cf2555aff7733b97b99}. The
\passthrough{\lstinline!omega\_symbol\_game!} crate tree at that commit
and at the code baseline is
\nolinkurl{75d77aa5e1f30dfc102f00c8f7cb6159418dec63}. This is retained
historical execution evidence. The 2026-10-02 package-local receipt does
not record a fresh run of these native targets.

The report is in this deposit at
\nolinkurl{docs/historical_native_run/report_second_revision_2026-09-28.md}.
The two JSON files beside it record the package review pipeline from the
same window. They do not name these four Cargo targets. Raw Cargo
standard output was not retained. The public source copy is
\nolinkurl{audited_source/pfp-iv-native-source-52cb5457.tar.gz}, with
member digests in \nolinkurl{audited_source/SOURCE_MEMBERS.sha256}.
Reproduction uses \passthrough{\lstinline!--locked!} and that pinned
tree, including its workspace and \passthrough{\lstinline!Cargo.lock!}.
A later checkout is a different tree.

For Python, collect and run
\passthrough{\lstinline!test\_packet\_contracts.py!},
\nolinkurl{test_exact_functional_port.py},
\nolinkurl{test_functional_boundary.py}, and
\nolinkurl{test_ffi_seam_isolation.py} under
\nolinkurl{omega_engine_symbol_game/tests/}. The exact-functional suite
contains scripted refusal cases and native cases. Its native cases can
skip when the executable is missing; a receipt must retain that
distinction. A scripted host tests adapter handling, not native
arithmetic.

A complete receipt records commands, exits, interpreter and Rust
versions, source digests, binary digest, environment-selected host, and
separate passed, failed, and skipped counts. Descriptive source
expectations in this paper must not be rewritten into measured results
until that receipt exists. Provider calls and model downloads are
unnecessary for these checks.

The package itself needs no repository checkout. From the package
directory, run one test file at a time with
\passthrough{\lstinline!py -3.12 -m pytest code/tests/<file>!}, build
the Lean modules with \passthrough{\lstinline!lake build!} in
\passthrough{\lstinline!lean4/!}, and replay the axiom census with
\passthrough{\lstinline!py -3.12 tools/replay\_lean\_axiom\_census.py --replay .!}.
The package tests check its legacy companion modules and the
weighted-cycle note. They run the governor's exact Python enumeration
unless a native graph executable is configured, so they are not evidence
for the native host described above.

\subsection{10. Limitations and practical
use}\label{limitations-and-practical-use}

PFP's useful guarantee is local: a consumer can inspect which finite
calculation was requested, which checks were replayed, which identities
were compared, and why admission remained unavailable. The inspected
host does not provide carrier-changing transport, arbitrary licensed
lifts, multiple candidates, or a fully authenticated deployment path.

The architecture also contains legacy helpers with weaker verification
contracts. Documentation must name the exact API rather than infer its
properties from shared terminology. Uniform RFC 8785 compliance, global
Rust/Python parity, and complete source-to-binary custody remain outside
the demonstrated scope. The companion code in this package is one
example. Its governor computes the graph minimum by exact Python
enumeration when no native executable is configured. Until this revision
it also substituted that enumeration whenever a configured native call
raised, including when the native adapter refused an input; it now lets
every native refusal or error propagate and refuses floating-point or
Boolean edge weights on both paths.

This separation enables honest use. A developer can reproduce a finite
experiment and retain its refusal without claiming that all proof
obligations or operational decisions have been settled. Future work
should target one missing boundary at a time, with a falsifying example
and a new source-bound receipt.

This paper assigns no authorization and changes no jurisdiction. It
carries no operational authority or institutional standing.

\subsection{Package verification}\label{package-verification}

The package verifier (\passthrough{\lstinline!verify\_deposit.py!})
checks the manifest and the two-way file census, and
\passthrough{\lstinline!tools/rebuild\_release.py --check!} checks the
release archive against the manifests. The current package-local record
is \nolinkurl{docs/release_validation_2026_10_02_review_repairs.json}.
It lists each command, tool version, exit status, and tested scope for
the checks run with this repair. It does not record a Cargo run or a
Lean build. \nolinkurl{docs/release_validation_2026_10_01.json} and
\nolinkurl{docs/release_validation_2026_10_02.json} keep their original
dates and scope.

\subsection{References}\label{references}

\begin{itemize}
\tightlist
\item
  Necula, G. C. (1997).
  \href{https://doi.org/10.1145/263699.263712}{Proof-carrying code}. In
  \emph{Proceedings of the 24th ACM SIGPLAN-SIGACT Symposium on
  Principles of Programming Languages (POPL '97)}, pp.~106--119. ACM.
  DOI: 10.1145/263699.263712.
\item
  Rundgren, A., Jordan, B., and Erdtman, S. (2020).
  \href{https://doi.org/10.17487/RFC8785}{JSON Canonicalization Scheme
  (JCS)}. RFC 8785 (Informational), RFC Editor, June 2020. DOI:
  10.17487/RFC8785.
\end{itemize}

Implementation links throughout the paper identify the exact audited
source revision. Section 9 identifies the historical native run and its
limits. The public inspection copy is
\nolinkurl{audited_source/pfp-iv-native-source-52cb5457.tar.gz}.

\end{document}
